\chapter{Aim and Scope}

\section{Purpose of this blueprint}

This blueprint records the first formalization target for the repository:
the narrow corridor from \IUT{} III, Theorem 3.11 to Corollary 3.12.
It is not a blueprint for the whole proof of the \(abc\) conjecture, and it
does not treat the correctness of \IUT{} as an assumption.  Its job is to make
the objects, comparison levels, indeterminacies, and real-valued estimates
visible enough that \leanfour{} can distinguish proved consequences from
source obligations.

\begin{definition}[Stage 1 target]
  \label{def:stage1-target}
  \lean{Iut.Stage1.theorem311ToCorollary312Labels}
  \leanok
  The first target is the implication
  \[
    \text{\IUT{} III, Theorem 3.11}
      \Longrightarrow
    \text{\IUT{} III, Corollary 3.12}.
  \]
  Following Mochizuki's 2026 formalization report, we refine this into the
  route
  \[
    \text{Theorem 3.11}
      \xrightarrow{\APT+\IPL}
    \text{``Theorem 3.11.5''}
      \xrightarrow{\SHE+\IPL}
    \text{Corollary 3.12}.
  \]
  The intermediate label ``3.11.5'' is not a separate theorem in the source
  papers.  It names the reorganized boundary obtained by moving the
  holomorphic hull plus determinant operation before the final simultaneous
  comparison of the \ThetaPilot{} and the \qPilot{}.
\end{definition}

\begin{remark}
  \label{rem:neutrality}
  This blueprint is neutral about the Mochizuki--Scholze--Stix dispute.  A
  completed Stage 1 route would say that the stated source hypotheses imply
  the Corollary 3.12 inequality.  It would not by itself prove that those
  source hypotheses are constructible from \IUT{} I--III.
\end{remark}

\section{Current repository status}

\begin{proposition}[Three levels of availability]
  \label{prop:availability-levels}
  The current Lean development distinguishes three kinds of availability.
  First, elementary ordered-real consequences are proved directly.  Second,
  route theorems prove that named source-facing records imply the desired
  signed comparison.  Third, some records are still obligations standing for
  Hodge-theater, Frobenioid, log-Kummer, hull, determinant, \SHE{}, \IPL{},
  and \APT{} constructions not yet built from the source papers.
\end{proposition}

\begin{proof}
  The code contains no proof placeholders in the current source files, and the
  project builds through the experiment surface.  The remaining gap is not a
  Lean elaboration gap; it is the mathematical construction of the source
  records consumed by the route theorems.
\end{proof}

\begin{definition}[Read order]
  \label{def:read-order}
  A reader should use the following order.  First read the source-facing
  statement of the Stage 1 route in \cref{ch:stage1-source}.  Then read the
  formal interface in \cref{ch:formal-interface}.  The actual conditional
  route through the corollary is in \cref{ch:corollary-route}.  The unresolved
  source obligations are collected in \cref{ch:source-obligations}.  Later
  consequences are only sketched in \cref{ch:after-corollary}.
\end{definition}

\begin{definition}[Two status axes]
  \label{def:two-status-axes}
  Each blueprint node should be read with two independent status axes.  Its
  Lean status says whether the displayed declaration checks, proves a theorem,
  or only records a diagnostic report.  Its source status says whether the
  corresponding mathematics has been constructed from \IUT{} I--III, is a
  conditional record field, is still a source obligation, or is only a later
  roadmap node.  In this project, a Lean-checked record can still be a
  source-level obligation.
\end{definition}
